\section{Discussion}
\label{sec:discussione}

\subsection{Interpretation of the results}
\label{sec:discussione-risultati}

Agreement between the solver and the zero-dimensional programs supports the
correct integration of the model into OpenFOAM. Differences remain within
0.3\% for the cases evaluated with both approaches, including chemical
reactions, except for high-temperature nitrogen. Comparison with the reference
study further assesses the thermochemical model. In non-reacting cases, final
temperatures are consistent with energy conservation and agree with the reference values to
within 0.3\%, which is within the uncertainty associated with extracting data
from the figures. In the reacting cases, the final temperatures agree to
within 0.8\%, the nitrogen-mixture composition to within 2.5\% and the air
composition to within 15\%. Transient temperature differences remain within
a few percent for pure nitrogen, the nitrogen-oxygen mixture and air. The
largest temperature discrepancy occurs in the vibrational temperature of the
molecular-atomic nitrogen mixtures, with deviations of 10--20\%.

\subsection{Sources of the discrepancies}
\label{sec:discussione-origine}

The delayed vibrational heating in molecular-atomic nitrogen is most plausibly
associated with the averaging of relaxation times over collision partners in
\mpp, which differs from the reference study. The close agreement for pure
nitrogen supports this interpretation, since averaging affects only mixtures
with multiple collision partners. The equilibrium state is unaffected by the
relaxation-time formulation, as it is determined by the conserved energy and
the thermodynamic model. This explanation remains a hypothesis; testing it
requires implementing the averaging procedure used in the reference study.
For pure nitrogen and the nitrogen-oxygen mixture, temporal shifts of a few
percent remain, and their origin has not been isolated.

\subsection{Methodological implications}
\label{sec:discussione-lezioni}

Reproducing the reference results required reconstructing modelling choices
that the reference study leaves unspecified or describes only partially.
These include the treatment of electronic energy, the energy used in the
relaxation source term, the exponent in Park's dissociation temperature and
the reaction-rate constants for air. The sensitivity comparisons in
Section~\ref{sec:risultati} show that these choices can produce deviations
larger than the residual differences: up to 56\% in vibrational temperature,
a factor of 3 in air species densities and approximately 25\% in final
temperature. The reference study also illustrates this sensitivity by
attributing its difference from LeMANS for molecular-atomic nitrogen to a
different convention in Park's correction.

Two complementary checks guided the selection of these modelling assumptions.
First, the equilibrium temperature obtained from energy conservation is
independent of the relaxation model and provides a check on the thermodynamic
data and energy conventions. Second, the curves extracted from the vector
figures enable a quantitative comparison with the reference study. Because the
solver and zero-dimensional programs share the source-term functions, their
agreement alone cannot exclude errors in those functions; the reference
curves and equilibrium checks provide complementary
evidence.
